\documentclass[11pt]{amsart}
\usepackage{geometry}                % See geometry.pdf to learn the layout options. There are lots.
\geometry{letterpaper}                   % ... or a4paper or a5paper or ... 
%\geometry{landscape}                % Activate for for rotated page geometry
%\usepackage[parfill]{parskip}    % Activate to begin paragraphs with an empty line rather than an indent
\usepackage{graphicx}
\usepackage{amssymb}
\usepackage{epstopdf}
\DeclareGraphicsRule{.tif}{png}{.png}{`convert #1 `dirname #1`/`basename #1 .tif`.png}

\title{PS 2.15}
%\date{}                                           % Activate to display a given date or no date

\begin{document}
\maketitle
%\section{}
%\subsection{}
Construya un diagrama de flujo que permita realizar operaciones aritméticas elementales, según la clave ingresada.
\begin{table}[htbp]
  \centering
  \begin{tabular}{|c|c|}
    \hline
    \multicolumn{2}{ |c| }{\textbf {Tabla 6.2}} \\ \hline
    Clave & Operación \\     \hline
    \hline
    + & 5UMA \\ \hline
    - & RE5TA \\ \hline
    * & MULTIPLICACIÓN \\ \hline
    / & DIVISIÓN \\ 
    \hline
  \end{tabular}
\end{table}
Imprima la clave ingresada y el resultado de la operación. \newline 
Datos:  \hspace{0.5cm}  OPER1, OPER2, CLAVE \newline 
Donde: 

 \setlength{\parindent}{1cm}
OPER1 \hspace{0.5cm}  es una variable de tipo real que representa el primer operando. 

 \setlength{\parindent}{1cm}
OPER2 \hspace{0.5cm}  es una variable de tipo real que expresa el segundo operando.

 \setlength{\parindent}{1cm}
CLAVE \hspace{0.5cm}  es una variable de tipo caracter que representa el tipo de operación aritmética que se va a realizar.
\end{document}  